\documentclass{article}
\usepackage{amsmath,hyperref}
\begin{document}
\section*{A task-relative representation of the named receipt}
\paragraph{Comparison and scope.}
Keep the native Cave Circuit registration and dispatch paths of PR 859,
the same receiver and debits, the same separate preparation, and the
same exact SQLite/RWARE settlement and restart check. Vary two things
independently: full versus compact suffix, and a dedicated early cue
versus an early byte of the selected encoding. Both suppliers pay for
their early send attempts. The sender/scope slot, codec identity and
offset are fixed by the declared comparison envelope; changing or
negotiating them is not free. This is a held-in controlled replay, not a
field measurement, independent mission or architectural acceptance.

\paragraph{Exact reconstruction, not invention.}
The pinned warehouse station writes the same encoded outcome $v$ into
the log and the named receipt under one transaction. Keep every
receipt's presence, intent, payload, and tick. Recover only $v$ from
the unique log row with that tick and payload. The decoder cannot use
the log to create a receipt that was absent. At encode time, require
that the join is unique and recovers every original value. For a fixed
result lookup $F$, the list of projected receipt keys can be expanded
to exactly the original list if and only if $F$ agrees at every
original key. The Lean file checks both directions and the empty case.
The proof is about the value projection; uniqueness and lookup
eligibility are code-side and source obligations, not granted by Lean.
An archive outside this schema gets no compact encoding in the test;
an unrestricted system would have to identify a full fallback or defer.

After decoding, require the same canonical full candidate on clean input,
and submit it to the unchanged digest and staged SQLite verification;
successful installation and an exact keyed reply are stronger than
decoding. A missing join has no candidate. A wrong key or a newly
consistent but wrong result can decode and still fail verification.
The experiment tracks delivery, decode, and installation separately.

\paragraph{Measured comparison and residual.}
The compact suffix is 165 bytes instead of 173 for the named archived
result. The P sender's old charge is $d+114$, and both early formats
spend one further byte. The independently varied four limits on $d$
are therefore 32, 33, 40, and 41 for full/dedicated, full/data-bearing,
compact/dedicated and compact/data-bearing respectively. The Q sender
can move its second diagnostic from tick 10 to tick 7 with either
compact format. With the full format, the same early diagnostic pushes
required capacity over 320. Both diagnostics and independent warehouse
work remain counted. The strongest equally informed direct scheduler
with the same codec ties. The earlier advantage against a full
fixed-length dedicated cue is not a comparison against task-relative
coding. Neither this projection nor a successful round trip certifies a
shortest possible encoding. Arithmetic here counts the tested slot and
already-agreed schema only, not any new metadata, distribution or
validation work in another interface.

A route that disappears after the early callback can still sink the
charged suffix in either format. Enlarging the accepted continuation
family to allow the incompatible late near registration makes both
formats withhold under the existing selector. Shorter representation
removes a capacity obstruction, not a route, copy, or authority
obstruction. No retained payload or successful check is credited as a
new authority to act.
\end{document}
